
\makeatletter
\let\@oldaddcontentsline\addcontentsline
\renewcommand{\addcontentsline}[3]{}
\section*{附录}
\let\addcontentsline\@oldaddcontentsline
\makeatother

\subsection*{A. 数据采集与清洗代码片段}

数据标准化处理示例代码（Python）：

\begin{verbatim}
import pandas as pd
import numpy as np
from sklearn.preprocessing import StandardScaler

# 加载原始数据
df = pd.read_csv('raw_data.csv')

# 缺失值处理
df.fillna(df.mean(), inplace=True)

# 异常值检测与处理
Q1 = df.quantile(0.25)
Q3 = df.quantile(0.75)
IQR = Q3 - Q1
mask = ~((df < (Q1 - 1.5 * IQR)) | (df > (Q3 + 1.5 * IQR)))
df = df[mask].dropna()

# 数据标准化
scaler = StandardScaler()
df_scaled = pd.DataFrame(
    scaler.fit_transform(df),
    columns=df.columns
)
\end{verbatim}

\subsection*{B. 主要技术栈}

\begin{table}[H]
  \centering
  \caption{系统开发技术栈}
  \label{tab:tech_stack}
  \begin{tabularx}{\textwidth}{LLL}
    \toprule
    \textbf{功能模块} & \textbf{开发语言/框架} & \textbf{关键库/工具} \\
    \midrule
    数据处理层 & Python & Pandas, NumPy, Scikit-learn \\
    数据库 & PostgreSQL, MongoDB & SQLAlchemy, PyMongo \\
    后端API & Python & Flask, FastAPI, RESTful \\
    前端Web & JavaScript & React, Vue.js, D3.js, Echarts \\
    移动应用 & React Native / iOS/Android & Flutter \\
    部署运维 & Linux & Docker, Kubernetes, Jenkins \\
    \bottomrule
  \end{tabularx}
\end{table}

\subsection*{C. 系统架构图}

\begin{table}[H]
  \centering
  \caption{系统架构及模块说明}
  \label{tab:architecture}
  \renewcommand{\arraystretch}{1.2}
  \begin{tabularx}{\textwidth}{LXL}
    \toprule
    \textbf{架构层} & \textbf{主要模块} & \textbf{功能描述} \\
    \midrule
    数据接入层 & 多源数据适配器、ETL流程 & 接收来自HIS、LIS、PACS等系统的数据 \\
    数据存储层 & 关系数据库、文档数据库、缓存系统 & 存储患者信息、检查结果、操作日志等 \\
    数据处理层 & 清洗、转换、聚合、计算引擎 & 数据标准化、特征工程、指标计算 \\
    分析计算层 & 统计分析、机器学习、深度学习模块 & 健康评估、风险预测、方案推荐 \\
    应用业务层 & 医生工作站、患者门户、管理后台 & 提供医患交互的业务功能 \\
    展现层 & Web界面、移动App、大屏展示 & 数据可视化和用户交互 \\
    \bottomrule
  \end{tabularx}
\end{table}

\subsection*{D. 患者隐私保护措施}

\begin{itemize}[itemindent=2em]
\item \textbf{数据加密：} 采用AES-256算法对敏感数据进行加密存储，TLS 1.3对数据传输进行加密。

\item \textbf{访问控制：} 实现基于角色的访问控制(RBAC)和属性的访问控制(ABAC)，确保数据只能被授权人员访问。

\item \textbf{审计日志：} 对所有关键操作进行记录和监控，定期审计可疑行为。

\item \textbf{数据脱敏：} 在测试环境使用脱敏数据，防止个人信息泄露。

\item \textbf{隐私政策：} 建立透明的隐私政策，获得患者明确同意。

\item \textbf{法规合规：} 严格遵守《个人信息保护法》、《数据安全法》等相关法律法规。
\end{itemize}

\subsection*{E. 性能测试报告摘要}

\begin{table}[H]
  \centering
  \caption{系统性能指标测试结果}
  \label{tab:performance}
  \begin{tabularx}{\textwidth}{CCCC}
    \toprule
    性能指标 & 测试条件 & 实测结果 & 指标要求 \\
    \midrule
    平均响应时间 & 100并发用户 & 328ms & <500ms \\
    最大吞吐量 & 1000并发 & 8500req/s & >5000req/s \\
    数据库查询速度 & 百万级数据 & 245ms & <1000ms \\
    CPU占用率 & 正常运行 & 45\% & <70\% \\
    内存占用 & 长期运行 & 2.3GB & <4GB \\
    可用性 & 30天运行 & 99.98\% & >99.9\% \\
    \bottomrule
  \end{tabularx}
\end{table}

\subsection*{F. 项目文件清单}

\begin{table}[H]
  \centering
  \caption{项目主要文件及说明}
  \label{tab:appendix}
  \renewcommand{\arraystretch}{1.2}
  \begin{tabularx}{\textwidth}{LX}
    \toprule
    \textbf{文件名} & \textbf{说明} \\
    \midrule
    data\_cleaning.py & 患者数据清洗与标准化处理代码 \\
    health\_score\_model.py & 健康评估模型实现代码 \\
    risk\_prediction\_model.py & 疾病风险预测模型（支持多种算法） \\
    recommendation\_engine.py & 个性化方案推荐系统 \\
    visualization.py & 数据可视化相关函数 \\
    database\_schema.sql & 数据库表结构定义 \\
    requirements.txt & Python依赖包列表 \\
    README.md & 系统安装与使用说明 \\
    试验数据集 & 系统验证用的脱敏患者数据 \\
    测试报告 & 单元测试、集成测试、性能测试报告 \\
    \bottomrule
  \end{tabularx}
\end{table}
